% Appendix B: material compressed out of Section 4 on Day 4 (2026-09-12). Verbatim.
\section{Results: full detail}
\label{app:results}

\subsection{E12 and the out-of-scope externals}
\label{app:e12}
\textbf{External, scope condition not met.} Two POPGym game environments
were tested and neither shows the effect. On BattleshipEasy the
capability gates failed outright (memoryless and memory arms both at the
floor, \(-0.441\) and \(-0.446\)), so P12.1--P12.3 are \emph{not
scored}. On MineSweeperMedium the gates passed convincingly --- our
memoryless arm reproduces POPGym's published MLP score to 0.012, and the
memory advantage at \(\lambda = 1.0\) is \(+0.121\) --- and all three
registered predictions then failed: the measured advantage at the
community default \(\lambda = 0.95\) is \(+0.091\), \emph{smaller} than
at \(\lambda = 1.0\), the opposite of the registered direction (P15.1),
and no threshold structure appears at the computed knee (P15.3).

We registered these as positive predictions and they failed as
registered. The reading, which the scope law itself supplies: in both
environments the latent state is fixed at reset and actions reveal
information about it rather than steering it, so the mechanism has no
equilibrium to displace. That distinction --- actions steering the
latent versus actions revealing it --- was not in our claim set when E12
was designed, and the $\tau$ used to place the knee was an episode length
rather than a relaxation time of a mean-reverting coordinate. Both are
our errors, and they are why P15 is reported as three failures rather
than reinterpreted. What survives is a boundary with evidence on both
sides of it: the effect appears where the theory says it should
(P9.1--P9.3) and is absent where the theory says it should be (G2,
P15.1--P15.3), on external benchmarks, under the same harness.

Independently, the released diagnostic (\S5.6) finds no trace of the
mechanism's signature in E12's training curves: memoryless-arm
excursions are at noise level and flat across \(\lambda\) (1, 0, 1
flagged runs of 5 below the knee against 2, 1, 1, 0, 1 above), rather
than concentrated below it. The instrument reports the effect absent,
not swamped.


\subsection{Mechanism misses and the E10 drift decomposition}
\label{app:mechanism}\label{app:drift}
Two observations bound what the demotion means, and neither is offered
as a rescue. The three P11.4 misses are the three scored cells with the
smallest positive \(S\) (\(+2.3\), \(+6.0\) and
\(+4.4 \times 10^{-4}\)), each adjacent to its instance's knee; and
resolvability is asymmetric at this budget --- growth resolved down to
\(|S| = 1.3\times10^{-4}\) while decay resolved only at
\(S \geq 1.1\times10^{-2}\), with perfect separation between the groups.
The registered floor was calibrated on a growth cell and applied
symmetrically, which is a design error we would not repeat. The minimal
reading of the three misses is ``unresolved at budget''; a genuine sign
error at \(|S| < 10^{-3}\) cannot be excluded on these data.

Separately, P11.1 and P11.2 failed: 19 of 32 cells drifted outside the
15\% gain tolerance. The decomposition is verified and reported in \S6
--- all 96 runs drifted downward (systematic, not diffusive), the
analytic cost map holds to \(\leq 1.05\%\) at every measured gain
including all failing cells, and failure aligns with the branch of the
equilibrium map rather than with \(\lambda\). P11.3 is estimated from
the shape of that intact cost profile, which is why it survives what
P11.1 and P11.2 did not.


